\documentclass[11pt,a4paper]{article}
\usepackage[margin=22mm]{geometry}
\usepackage{amsmath,amssymb,booktabs,array}
\usepackage[hidelinks]{hyperref}
\setlength{\parindent}{0pt}
\setlength{\parskip}{5pt}
\emergencystretch=3em
\raggedright
\begin{document}
{\Large Alignment ladder: exact local certificates and the remaining shared update}\\[4pt]
8 October 2026 --- methodology review; no traffic simulation or model adoption

\textbf{Conclusion.} The unchanged uniform field fails a local compatibility
test. Restriction-aware and continuous-wave descriptions have exact local
constructions, including own motion. Neither a completed numerical update nor
an applicable full continuum result has yet been supplied for all approved
lane-transfer, node, record and reset interactions. This is not a theorem that
the model cannot be constructed. A full arbitrary-data continuum theorem is
not required if an admissible, coherent finite numerical model is qualified.

\section*{1. Quantities and the unchanged projection obstruction}
Traffic density $\rho$ and jam density $\kappa$ have units equivalents/m.
Free speed $u$, backward-wave magnitude $w$ and restriction speed $s$ have
units m/s. Road flux $f$ has units equivalents/s. The triangular diagram is
\begin{equation}
 f(\rho)=\min\{u\rho,w(\kappa-\rho)\},\qquad
 \rho_c=\frac{w\kappa}{u+w},\qquad Q=u\rho_c.
\end{equation}
Let $C(t)$ be a local cumulative incoming count (with a fixed gauge offset),
and $N(t,x)$ the count remaining after subtracting traffic upstream of $x$:
\begin{equation}
 N(t,x)=C(t)-\int_0^x\rho(t,\xi)\,d\xi,\qquad
 N_x=-\rho,\quad N_t=f(\rho).
\end{equation}
The last identity is for a source-free segment with its correctly accounted
upstream port. A restriction position $z_j(t)$ has count $h_j=N(t,z_j(t))$.
Its speed $v_j=\dot z_j$ gives
\begin{equation}
 \dot h_j=f(\rho)-v_j\rho=0
\end{equation}
on the approved zero-relative-passing path. Records are effect metadata;
they add no population to these integrals.

Recheck the earlier 30 m cell, $\Delta t=1$, $u=30$, $w=5$, $\kappa=3/20$,
initial density $1/100$, free exit, sufficient input, and coincident old
records with $h_1=h_2=0$, desired speeds 5 and 15. The own fan at one second
has $\rho=3/40$ on $[0,5]$ and vacuum on $(5,30)$; positions are 5 and 15.
The initial $3/10$ mass exits and $3/8$ enters. Consequently
\begin{equation}
 m(1)=C(1)=\frac38,\qquad N(1,5)=N(1,15)=N(1,30)=0.
\end{equation}
Uniform projection instead uses $\bar\rho=(3/8)/30=1/80$ and yields
\begin{equation}
 \bar N(1,5)=\frac5{16},\qquad \bar N(1,15)=\frac3{16},\qquad
 \int_5^{15}\bar\rho\,dx=\frac18.
\end{equation}
Thus positions, old labels and uniform reconstruction cannot all hold.
The parent CFL equality $u\Delta t=30$ is satisfied. The example is local,
not a reached metering witness or a proof against every CTM realization.

\newpage
\section*{2. A new exact two-restriction own-field evolution}
Use the same diagram in $[0,30]$. At time zero the field consists of two
congested blocks separated by vacuum. Let
\begin{equation}
 r_1=\frac{w\kappa}{w+5}=\frac3{40},\qquad
 r_2=\frac{w\kappa}{w+10}=\frac1{20}.
\end{equation}
For $0\le t\le1/2$, prescribe only the external upstream demand needed to
maintain the own admissible inflow $q_0=5r_1=3/8$; exit is free. The explicit
solution, with old desired speeds 5 and 10, is
\begin{equation}
 \rho(t,x)=\begin{cases}
 r_1,&0\le x<5+5t,\\
 0,&5+5t<x<15+10t,\\
 r_2,&15+10t<x<20+10t,\\
 0,&20+10t<x\le30,
 \end{cases}\qquad
 z_1(t)=5+5t,\quad z_2(t)=20+10t.
\end{equation}
There are no interior lateral sources, encounters or exits in this local
window. Ordinary shocks at the rear of the second block and the nonclassical
zero-passing jumps at each restriction determine the same motion and field.
Ahead of each restriction is vacuum, so K.A permits its held desired speed;
no positive interior recovery is used. Within a block, $f(r_i)/r_i=s_i$.

The second block's left jump is an ordinary entropy shock: its speed is
$(f(r_2)-0)/(r_2-0)=10$, between the triangular diagram's slopes 30 and $-5$.
At each right jump,
\begin{equation}
 f(r_1)-5r_1=0,\qquad f(r_2)-10r_2=0.
\end{equation}
Take gauge $C(0)=m(0)=5r_1+5r_2=5/8$. Then
\begin{align}
 m(t)&=(5+5t)r_1+5r_2=\frac58+\frac38t=C(t),\\
 h_1(t)&=C(t)-(5+5t)r_1=\frac14,\\
 h_2(t)&=C(t)-m(t)=0.
\end{align}
At $t=1/2$, positions are $15/2$ and 25, the second block begins at 20,
and $m=C=13/16$. The intervening mass between the records is exactly $1/4$.
All flow, storage, position and count calculations agree in one own field.

\textbf{Scope.} This is an analytic local construction, not a public
E1/metering trajectory or the complete sourced update. It proves more than
reset interpolation, but not the full task. The exact field has two downward
density jumps and cannot be represented by a single internal discontinuity.
That does not exclude extending a subcell method to multiple pieces, or a
qualified numerical approximation that preserves all required anchors.

\newpage
\section*{3. What a genuine bridge must advance}
For moving interval edges $a(t)$ and $b(t)$, conservation gives the Reynolds
balance
\begin{equation}
 \frac{d}{dt}\int_{a(t)}^{b(t)}\rho\,dx
 = [f-\dot a\rho]_{a(t)}-[f-\dot b\rho]_{b(t)}.
\end{equation}
At an active zero-passing record boundary, its relative flux is zero. At a
fixed transfer port, the flux comes from the same accepted node/source solve,
not a second budget. This balance determines mass \emph{after the boundary
fluxes and paths have been solved}; it is not itself that solve.

For completed compatible anchors $(x_i,N_i)$, ordered by position,
\begin{equation}
 0\le N_i-N_{i+1}\le\kappa(x_{i+1}-x_i),\qquad
 \rho_i=\frac{N_i-N_{i+1}}{x_{i+1}-x_i}.
\end{equation}
Coincident positions must have equal count. Distinct kinetic records remain.
The interpolation conserves counts and preserves vacuum when adjacent counts
are equal. It does not prescribe the next moving flux or its convergence.

Source IT uses a regular-step sending mass $S_a=\Delta t D_a$, lane-choice
rate $\pi_{ab}$ (s$^{-1}$), and dimensionless oriented fraction
\begin{equation}
 \theta_{ab}=\pi_{ab}\Delta t\quad(b\ne a),\qquad
 \theta_{aa}=1-\sum_{b\ne a}\theta_{ab},\qquad
 d_{ab}=\theta_{ab}S_a.
\end{equation}
The geometry advances $(i,a)$ to $(i+1,b)$ once. The source sampling, sending
and available receiving define a finite numerical step. Replacing $D_a$ by
an instantaneous boundary trace while keeping the same fractions is an
explicit attempted adaptation, not a proved equivalence. Internal numerical
steps must not silently resample $\theta$ or E1/positive recovery.

For accepted transfer mass $F_{ab}$, paired balances give
\begin{equation}
 m_a^+=m_a+\sum_bF_{ba}-\sum_bF_{ab}.
\end{equation}
After including outside waiting, external offers and exits, internal transfers
cancel. This is conditional conservation: it does not establish the own
admissibility or correctness of every $F_{ab}$.

K.A uses a held desired speed and the running minimum of encountered own
traffic caps within a public interval, with position obtained by integration.
Ending speed is retained separately. E1 attaches records to the completed
own count at the next public clock. A complete numerical bridge must specify
the moving/source/port solve preceding those operations. A continuous count
formula supplied with external paths does not replace this requirement.

\newpage
\section*{4. A direct continuous proportional-merge transplant fails consistency}
This is an independent local junction diagnostic. It is \emph{not} a reached
oriented-IT/metering case and does not refute native CTM. It identifies why
the step from a finite source rule to a boundary-trace rule needs qualification.

Let $d_1=1/5$, $d_2=2/5$ be demands on two independent incoming roads,
$R=3/10$ the single receiving supply, and both lane capacities
$Q=9/14$. Direct demand-proportional allocation gives
\begin{equation}
 g_i=\frac{d_i}{d_1+d_2}\min\{d_1+d_2,R\},\qquad
 (g_1,g_2)=\left(\frac1{10},\frac15\right).
\end{equation}
Both roads are restricted: $g_i<d_i$. Their initial free densities are
$\rho_i=d_i/u$. The admissible boundary traces for the restricted flows are
congested,
\begin{equation}
 \rho_i^*=\kappa-\frac{g_i}{w},\qquad
 \rho_1^*=\frac{13}{100},\quad \rho_2^*=\frac{11}{100},\quad
 D(\rho_1^*)=D(\rho_2^*)=Q.
\end{equation}
Each connecting shock has negative speed because
\begin{equation}
 \sigma_i=\frac{g_i-d_i}{\rho_i^*-\rho_i}<0.
\end{equation}
These are appropriate upstream congestion waves. Reapplying the direct node
formula to the states it generated instead returns
\begin{equation}
 (g_1^*,g_2^*)=\left(\frac3{20},\frac3{20}\right)
 \ne\left(\frac1{10},\frac15\right).
\end{equation}
Therefore this boundary rule is not invariant under its own Riemann states.
Continuous LTM literature requires a suitable invariant realization, and also
distinguishes CTM interior-state interpretations. The calculation does not
prove that such a source-equivalent completion is impossible.

\textbf{Do not attribute this failure to N.A.} With fixed priority $p=1/2$,
the selected work-conserving two-input rule is
\begin{equation}
 g_1=\min\{d_1,\max\{pR,R-d_2\}\},\qquad
 g_2=\min\{d_2,R-g_1\}.
\end{equation}
For the same demands it gives $(3/20,3/20)$, and the restricted traces with
demands $(Q,Q)$ give that same split. Another case $d_1=1/10$, $d_2=2/5$
gives $(1/10,1/5)$, unchanged after only restricted road 2 raises its demand
to $Q$. These checks are not a proof for all multi-input IT/moving-path nodes.

\newpage
\section*{5. Qualification, not a fabricated global obstruction}
\begin{center}
\begin{tabular}{>{\raggedright\arraybackslash}p{.18\textwidth}>{\raggedright\arraybackslash}p{.32\textwidth}>{\raggedright\arraybackslash}p{.38\textwidth}}
\toprule
Route & Established & Still required for full success\\
\midrule
Interface & Uniform projection contradiction & A richer own field, not only renamed labels\\
Adaptation & Anchor consistency and the explicit two-restriction wave evolution & Complete admissible moving/node/source numerical update for the accepted phases\\
Representation & Continuous local count/wave construction & Endogenous paths and source-equivalent transfer ports, not externally supplied histories\\
\bottomrule
\end{tabular}
\end{center}

An acceptable finite numerical model need not solve the continuum exactly.
It must nevertheless state its own source/flux update, conditions and how it
preserves the required metadata/field meaning. Ordinary field truncation error
does not justify an arbitrary mismatch at retained count boundaries.
Likewise a continuum method must specify how its paths and source histories
are obtained from its own field. Conservation alone does not do this.

\textbf{Source versus original calculation.} Sections 1--2 recheck/extend
the reviewer-derived local constructions under the previously selected
triangular diagram and zero passing. Section 4 is an original exact diagnostic
of a direct proportional-merge transplant, consistent with the source's
invariance distinction. It is not presented as a literature proof against
the entire approved IT graph. Section 3 records the source-time dependence
and conditional balances; no new physical source rule is adopted.

\textbf{Numerical alternatives.} A1 retains moving anchor boundaries. A2
retains internal discontinuities inside fixed cells. C supplies a continuous
count/wave reference. They can represent the same physics. Single-restriction
coupled schemes and existence results are method precedents, not automatically
the multiple-record, zero-passing, sampled-recovery/source task. A single-shock
source formula cannot reproduce the exact two-jump field above unchanged;
multi-piece extension remains possible.

\textbf{Final status.} No full success step; step 2 has the strongest local
positive. Step 1 is refuted only on its audited domain; steps 2--3 are
unestablished full constructions, not mathematical impossibility. The bounded
ladder is reviewed. Q.J.F / Stage 3 / continuation 67 remain open; Stage 4 is
unopened. Existing M.A, N.A, G1, E1 and K.A remain approved. Clearance, full
TST response, useful gradients and necessity of retained information are not
established by this companion.

\textbf{Supporting sources.} Full identities, access limits and extraction
are in the source audit. Main junction reference:
\href{https://arxiv.org/pdf/1405.7080}{Jin, continuous LTM, section 6}.
The proposed fixed-cell alternative uses
\href{https://ems.press/content/serial-article-files/36595}{Chalons et al., 2017}.
These links are evidence provenance, not replacement update laws.
\end{document}
